\documentclass[12pt,a4paper]{article}
\usepackage[utf8]{inputenc}

\usepackage{color}
%\usepackage[magyar]{babel}
\usepackage{listings}
\usepackage{graphicx}
\usepackage{enumerate}
\usepackage{tikz}
\usetikzlibrary{shapes,arrows}
\usetikzlibrary{positioning}

\renewcommand{\refname}{Irodalomjegyzék}
\renewcommand{\contentsname}{Tartalom}
\renewcommand{\lstlistingname}{Lista}
\usepackage[
	pdftitle={HA5KFU tanfolyam},
	pdfauthor={HA5KFU Rádióamatőr Klub},
	pdfsubject={HA5KFU tanfolyam},
	pdfcreator={latex},
	pdfkeywords={ },
	pdfpagemode=UseOutlines,
	pdfdisplaydoctitle=true,
	pdflang=hu,
	unicode
]{hyperref}

\pagestyle{plain}
\begin{document}
\begin{center}
\includegraphics[width=300pt,keepaspectratio]{figures/ha5kfu.eps}
\\[2cm]

{\huge \bfseries Forgalmazási ismeretek \\[2cm]}



\end{center}
\tableofcontents 
\newpage
\section{Jogszabályok}
\subsection{Rádióamatőr engedély}
\subsection{Hívójelek}

\newpage
\begin{thebibliography}{9}
\bibitem{5kdr_j} 
Budapest Fővárosi Rádióamatőr Klub 
\textit{Rádióamatőr vizsgára felkészítő jegyzet}. 
\\\texttt{https://www.ha5kdr.hu/}


\end{thebibliography}

\end{document}
